\chapter{Bipartite Matching}

In the previous chapter, we saw how assignment and transportation problems reduce directly to minimum-cost flow.
When capacities are all one and costs are uniform, the problem simplifies to pairing compatible items
across two disjoint sets so that no vertex is used more than once.

This is \emph{bipartite matching}. A worker takes one job, a student occupies one room, a processor
runs one task. Although the problem reduces to unit-capacity network flow, its combinatorial structure
yields direct algorithms: alternating paths that flip edges in place, and the Hopcroft--Karp batching
strategy that achieves $\Oh(E\sqrt{V})$ time without constructing a flow network.

\chapterline

\section{Matchings: the language}
\label{ch15:language}

\begin{definition}[Bipartite graph]
\label{ch15:def-bipartite}
A graph $\Gph=(V,E)$ is \emph{bipartite} if its vertices split into two disjoint
sets $L$ and $R$ (the left and right \emph{sides}, or \emph{color classes}) so
that every edge has one endpoint in $L$ and one in $R$. We write
$\Gph=(L\cup R,E)$ and set $n=\abs{V}=\abs{L}+\abs{R}$, $m=\abs{E}$.
\end{definition}

Equivalently, $\Gph$ is bipartite iff it has no odd cycle (a two-coloring of the
vertices exists iff every cycle has even length). A single breadth-first sweep
finds the bipartition or an odd cycle, so we will assume $L$ and $R$ are given.

\begin{definition}[Matching]
\label{ch15:def-matching}
A \emph{matching} $M\subseteq E$ is a set of edges no two of which share an
endpoint. A vertex incident to an edge of $M$ is \emph{matched} (or
\emph{saturated}); otherwise it is \emph{free} (or \emph{exposed}). A matching is
\emph{maximum} if no matching has more edges, and \emph{perfect} if every vertex
is matched (so $\abs{M}=\abs{L}=\abs{R}$).
\end{definition}

\begin{intuitionbox}
Picture the left column as people and the right column as chairs. A matching is a
collection of ``person sits in chair'' edges where no person sits in two chairs and
no chair holds two people. ``Maximum'' means: seat as many people as possible. The
question is never about one edge -- it is about the global tradeoff, because
greedily grabbing one edge can block two better ones.
\end{intuitionbox}

Here is the running example for the whole chapter. The left side is
$L=\set{a,b,c,d}$, the right side is $R=\set{1,2,3,4}$, and the edges are
\[
E=\set{a1,\;a2,\;b1,\;b3,\;c3,\;c4,\;d3}.
\]

\begin{figure}[H]
\centering
\begin{tikzpicture}[
  every node/.style={font=\small},
  L/.style={circle,draw=sourceblue,thick,fill=blue!8,minimum size=7mm},
  R/.style={circle,draw=targetgreen,thick,fill=green!8,minimum size=7mm},
  e/.style={edgegray}
]
\foreach \v/\y in {a/3,b/2,c/1,d/0} \node[L] (\v) at (0,\y) {$\v$};
\foreach \v/\y in {1/3,2/2,3/1,4/0} \node[R] (r\v) at (3.2,\y) {$\v$};
\draw[e] (a)--(r1); \draw[e] (a)--(r2);
\draw[e] (b)--(r1); \draw[e] (b)--(r3);
\draw[e] (c)--(r3); \draw[e] (c)--(r4);
\draw[e] (d)--(r3);
\node at (0,3.9) {\textbf{$L$}};
\node at (3.2,3.9) {\textbf{$R$}};
\end{tikzpicture}
\caption{The running bipartite graph $\Gph$. Left $L=\set{a,b,c,d}$, right
$R=\set{1,2,3,4}$.}
\label{ch15:fig-running}
\end{figure}

A greedy person scanning edges in the listed order might grab $a1$, then $b3$,
then $c4$, and stop: $d$ wants only chair $3$, which is taken, and $a$'s
alternative chair $2$ is now stranded with nobody left to ask cleanly. That
greedy run produces $\set{a1,b3,c4}$ of size $3$. But $\set{a2,b1,c4,d3}$ has
size $4$ and is perfect, so greedy fell one short. The rest of the chapter shows
how to repair such mistakes systematically.

\section{Augmenting paths and Berge's theorem}
\label{ch15:berge}

Fix a matching $M$. The key device is to look at paths that \emph{alternate}
between non-matching and matching edges.

\begin{definition}[Alternating and augmenting paths]
\label{ch15:def-aug}
Relative to $M$, a path is \emph{$M$-alternating} if its edges alternate between
$E\setminus M$ and $M$. An alternating path is \emph{$M$-augmenting} if it starts
and ends at \emph{free} vertices (necessarily one in $L$, one in $R$, and the
path has odd length).
\end{definition}

\begin{intuitionbox}
An augmenting path is a chain ``free person -- empty chair -- person -- their chair
-- person -- their chair -- $\dots$ -- empty chair''. Walk it and flip every edge:
the non-matching edges (the ``wants'') become matches, the matching edges become
free. Because the path starts and ends free and alternates, flipping turns $k$
matched edges into $k+1$ matched edges. Net gain: exactly one. It is a chain of
bumps: the first free person takes a chair, bumping its occupant, who takes the
next chair, bumping again, until the last bump lands on an empty chair.
\end{intuitionbox}

Formally, if $P$ is $M$-augmenting then the symmetric difference
\[
M' \;=\; M \oplus P \;=\; (M\setminus P)\cup(P\setminus M)
\]
is again a matching, and $\abs{M'}=\abs{M}+1$. The endpoints of $P$ go from free
to matched; every internal vertex keeps exactly one matched edge (its matched
edge swaps to the other side of the path), so $M'$ is a valid matching, and $P$
has one more non-matching edge than matching edge.

\begin{figure}[H]
\centering
\begin{tikzpicture}[
  every node/.style={font=\small},
  L/.style={circle,draw=sourceblue,thick,fill=blue!8,minimum size=7mm},
  R/.style={circle,draw=targetgreen,thick,fill=green!8,minimum size=7mm},
  m/.style={trapred,very thick},
  e/.style={edgegray}
]
\begin{scope}
\foreach \v/\y in {a/3,b/2,c/1,d/0} \node[L] (\v) at (0,\y) {$\v$};
\foreach \v/\y in {1/3,2/2,3/1,4/0} \node[R] (r\v) at (2.6,\y) {$\v$};
\draw[e] (a)--(r2); \draw[e] (b)--(r1); \draw[e] (c)--(r3);
\draw[m] (a)--(r1); \draw[m] (b)--(r3); \draw[m] (c)--(r4);
\draw[e,dashed] (d)--(r3);
\node at (0,3.9) {$L$}; \node at (2.6,3.9) {$R$};
\node at (1.3,-1.1) {\textbf{(a)} $M=\set{a1,b3,c4}$};
\end{scope}
\begin{scope}[xshift=7cm]
\foreach \v/\y in {a/3,b/2,c/1,d/0} \node[L] (\v) at (0,\y) {$\v$};
\foreach \v/\y in {1/3,2/2,3/1,4/0} \node[R] (r\v) at (2.6,\y) {$\v$};
\draw[m] (a)--(r2); \draw[m] (b)--(r1); \draw[m] (c)--(r4); \draw[m] (d)--(r3);
\draw[e] (a)--(r1); \draw[e] (c)--(r3); \draw[e] (b)--(r3);
\node at (0,3.9) {$L$}; \node at (2.6,3.9) {$R$};
\node at (1.3,-1.1) {\textbf{(b)} after augmenting along $P$};
\end{scope}
\end{tikzpicture}
\caption{Left: matching $M=\set{a1,b3,c4}$ (thick red) with free vertices
$d,2$ and the augmenting path $P=d\!-\!3\!-\!b\!-\!1\!-\!a\!-\!2$ (it alternates
$d3\notin M$, $3b\in M$, $b1\notin M$, $1a\in M$, $a2\notin M$). Right: flipping
$P$ yields the perfect matching $\set{a2,b1,c4,d3}$ of size $4$.}
\label{ch15:fig-augment}
\end{figure}

So an augmenting path is a certificate that $M$ is \emph{not} maximum. The harder
fact is the converse: if no augmenting path exists, $M$ already is maximum. This
is Berge's theorem, and it holds for general (not just bipartite) graphs.

\begin{theorem}[Berge, 1957]
\label{ch15:thm-berge}
A matching $M$ in a graph $\Gph$ is maximum if and only if $\Gph$ contains no
$M$-augmenting path.
\end{theorem}

\begin{whybox}
The forward direction is the easy half: if an augmenting path existed we could
flip it and get a strictly larger matching, so a maximum $M$ cannot admit one.
For the reverse direction, the key tool is the
symmetric difference of two matchings. Compare $M$ with a hypothetical
\emph{better} matching $M^\ast$. Their symmetric difference is forced into a very
rigid shape -- disjoint paths and even cycles -- and counting edges in those
pieces produces an augmenting path whenever $M^\ast$ is larger.
\end{whybox}

\begin{proof}
($\Rightarrow$) If $P$ is $M$-augmenting then $M\oplus P$ is a matching of size
$\abs{M}+1$, so $M$ is not maximum. Contrapositive: maximum $\Rightarrow$ no
augmenting path.

($\Leftarrow$) Suppose $M$ is not maximum; let $M^\ast$ be a matching with
$\abs{M^\ast}>\abs{M}$. Consider the edge set $H=M\oplus M^\ast$ and the subgraph
it induces. Every vertex $v$ is incident to at most one edge of $M$ and at most
one of $M^\ast$, hence to at most two edges of $H$. So in $H$ every vertex has
degree $\le 2$, which means each connected component of $H$ is either a simple
path or a simple cycle.

Within any component the edges alternate between $M$ and $M^\ast$: two
consecutive $H$-edges at a vertex cannot both be in $M$ (they would share that
vertex, impossible in a matching) nor both in $M^\ast$. Consequently every
\emph{cycle} component has even length and contains equally many $M$- and
$M^\ast$-edges. Since $\abs{M^\ast}>\abs{M}$ overall, some component must contain
strictly more $M^\ast$-edges than $M$-edges. That component cannot be a cycle, so
it is a path, and a path with one extra $M^\ast$-edge must begin and end with an
$M^\ast$-edge. Both endpoints of that path are therefore unmatched by $M$ (an
$M$-edge at an endpoint would extend the path), so the path is $M$-alternating
with free ends: an $M$-augmenting path. Contrapositive gives the claim.
\end{proof}

\begin{trapbox}
Two pitfalls. (1) Augmenting means flipping a path between \emph{free} endpoints;
flipping an alternating path that ends at a matched vertex does \emph{not}
increase the matching and can even break it. (2) Berge's theorem says a maximum
matching has \emph{no} augmenting path -- it does \emph{not} say a greedy/maximal
matching has none. A maximal matching (one you cannot extend by adding a single
edge) can be half the size of maximum; the running graph's greedy
$\set{a1,b3,c4}$ is maximal but not maximum precisely because an augmenting path
survives.
\end{trapbox}

Berge's theorem hands us an algorithm template immediately: start from the empty
matching and repeatedly find an augmenting path and flip it, until none remains.
Each flip adds one edge, so after at most $\min(\abs{L},\abs{R})$ flips we stop
with a maximum matching. Everything below is about finding those paths fast.

\section{The min-max duality: K\H{o}nig and Hall}
\label{ch15:konig}

Berge tells us when to stop. The dual theorems tell us \emph{why} the number we
stop at is what it is, by exhibiting a small obstruction whose size equals the
matching's. This is the same min-max flavor as max-flow/min-cut, and indeed it is
a special case of it.

\begin{definition}[Vertex cover]
\label{ch15:def-cover}
A \emph{vertex cover} is a set $C\subseteq V$ touching every edge: for each
$uv\in E$, $u\in C$ or $v\in C$. Write $\tau(\Gph)$ for the minimum cover size and
$\nu(\Gph)$ for the maximum matching size.
\end{definition}

For \emph{any} graph $\nu\le\tau$: a matching has $\nu$ disjoint edges, and a
cover must spend at least one distinct vertex per matched edge. K\H{o}nig's theorem
shows that for bipartite graphs, this inequality is always tight.

\begin{theorem}[K\H{o}nig, 1931]
\label{ch15:thm-konig}
In a bipartite graph, $\nu(\Gph)=\tau(\Gph)$: the maximum matching size equals
the minimum vertex cover size.
\end{theorem}

\begin{proof}
We have $\nu\le\tau$ always. For the reverse, let $M$ be a maximum matching;
Berge guarantees no $M$-augmenting path. Direct every non-matching edge from $L$
to $R$ and every matching edge from $R$ to $L$, and let $Z$ be the set of
vertices reachable by $M$-alternating paths that \emph{start at a free vertex of
$L$} (do a search from all free left vertices following the alternation rule).

Define the cover
\[
C \;=\; (L\setminus Z)\;\cup\;(R\cap Z).
\]
We claim $C$ is a vertex cover of size $\abs{M}$.

\emph{$C$ covers every edge.} Take any edge $uv$ with $u\in L$, $v\in R$.
\begin{itemize}
\item If $u\notin Z$ then $u\in L\setminus Z\subseteq C$, covered.
\item If $u\in Z$, then there is an alternating path from a free left vertex to
$u$. The edge $uv$ is either in $M$ or not. If $uv\notin M$, the search extends
the alternating path across $uv$ to $v$, so $v\in Z$, hence $v\in R\cap Z\subseteq
C$. If $uv\in M$, then $u$ is matched and therefore not a free starting vertex of
$Z$; since every directed edge entering $L$ in the alternating search is a matched
edge and $uv$ is the unique matched edge incident to $u$, the search could only
have reached $u$ by arriving \emph{along} $uv$ from $v$, so again $v\in Z$ and
$v\in C$.
\end{itemize}
Either way $uv$ is covered, so $C$ is a vertex cover.

\emph{$\abs{C}=\abs{M}$.} First, every vertex of $L\setminus Z$ is matched: free
left vertices are in $Z$ by definition, so a vertex outside $Z$ on the left is
not free. Second, every vertex of $R\cap Z$ is matched: a free right vertex in
$Z$ would be the end of an $M$-augmenting path, contradicting maximality. So
every vertex of $C$ is saturated by $M$. We now show each matched edge
contributes exactly one endpoint to $C$. Take a matched edge $u'v'$ with $u'\in
L,\,v'\in R$. If $v'\in R\cap Z$, then $v'$ was reached, and the only way the
alternating search enters the right side along a matched edge is from its left
endpoint, so the path reaches $v'$ via $u'$, forcing $u'\in Z$, i.e. $u'\notin
L\setminus Z$. Thus $u'$ and $v'$ are never both in $C$, and since every vertex
of $C$ is the endpoint of a matched edge, $\abs{C}=\abs{M}$. Combined with
$\tau\ge\nu=\abs{M}$ we get $\tau=\abs{M}=\nu$.
\end{proof}

Figure~\ref{ch15:fig-konig-cover} visualizes this four-quadrant decomposition: the alternating search partitions both $L$ and $R$ by reachability from free left vertices. The minimum vertex cover $C = (L\setminus Z)\cup(R\cap Z)$ takes precisely one endpoint of every matched edge, while no edge can cross from $L\cap Z$ to $R\setminus Z$ without expanding $Z$.

\begin{figure}[H]
\centering
\begin{tikzpicture}[>=Stealth, thick, scale=0.92,
  vnode/.style={circle, draw=nodeblue, fill=blue!12, minimum size=7.5mm, font=\small\bfseries},
  rnode/.style={circle, draw=targetgreen, fill=green!12, minimum size=7.5mm, font=\small\bfseries},
  coverbox/.style={draw=quotientorange, fill=orange!12, rounded corners=6pt, dashed, line width=1pt},
  normbox/.style={draw=edgegray!60, fill=gray!5, rounded corners=6pt, dashed}]

  % Quadrant 1: L \ Z (in Cover C)
  \draw[coverbox] (-0.4, 2.5) rectangle (1.8, 5.0);
  \node[anchor=north west, font=\scriptsize\bfseries, text=quotientorange] at (-0.3, 4.9) {$L\setminus Z \subseteq C$};
  \node[vnode] (u1) at (0.7, 3.5) {$u_1$};

  % Quadrant 2: L \cap Z (not in cover, contains free left vertices)
  \draw[normbox] (-0.4, -0.6) rectangle (1.8, 2.1);
  \node[anchor=north west, font=\scriptsize\bfseries, text=sourceblue] at (-0.3, 2.0) {$L\cap Z$};
  \node[vnode] (u2) at (0.7, 1.2) {$u_2$};
  \node[vnode] (ufree) at (0.7, 0.0) {$u_{\text{free}}$};

  % Quadrant 3: R \ Z (not in cover)
  \draw[normbox] (5.2, 2.5) rectangle (7.4, 5.0);
  \node[anchor=north east, font=\scriptsize\bfseries, text=targetgreen] at (7.3, 4.9) {$R\setminus Z$};
  \node[rnode] (v1) at (6.3, 3.5) {$v_1$};

  % Quadrant 4: R \cap Z (in Cover C)
  \draw[coverbox] (5.2, -0.6) rectangle (7.4, 2.1);
  \node[anchor=north east, font=\scriptsize\bfseries, text=quotientorange] at (7.3, 2.0) {$R\cap Z \subseteq C$};
  \node[rnode] (v2) at (6.3, 0.7) {$v_2$};

  % Matched edge 1: between L\Z and R\Z (covered by u1 in C)
  \draw[trapred, line width=2pt] (u1) -- node[midway, above=1pt, font=\scriptsize\bfseries, text=trapred] {$\in M$} (v1);

  % Matched edge 2: between L \cap Z and R \cap Z (covered by v2 in C)
  \draw[trapred, line width=2pt] (u2) -- node[midway, above=1pt, font=\scriptsize\bfseries, text=trapred] {$\in M$} (v2);

  % Unmatched edge from u_free to v2 (alternating search step)
  \draw[->, very thick, edgegray] (ufree) -- node[midway, below=1pt, font=\tiny, text=gray] {search $L\to R$} (v2);

  % Unmatched edge from u1 to v2 (both endpoints in cover)
  \draw[edgegray, dashed] (u1) -- (v2);

  % Forbidden cross: L \cap Z to R \setminus Z
  \draw[edgegray!40, dotted, line width=1pt] (u2) -- (v1);
  \node[draw=trapred, fill=white, rounded corners=2pt, font=\tiny\bfseries, text=trapred, inner sep=2pt] at (3.5, 2.3) {No edges: $L\cap Z \to R\setminus Z$};

  % Summary caption note
  \node[anchor=north, font=\scriptsize, align=center, text=gray!80!black] at (3.5, -1.0) {
    Cover $C = (L\setminus Z)\cup(R\cap Z)$ picks exactly one vertex per matched edge.\\
    Every non-matched edge is covered; no edges cross to $R\setminus Z$. $\abs{C} = \abs{M}$.
  };
\end{tikzpicture}
\caption{K\H{o}nig's minimum vertex cover construction. Alternating reachability from free left vertices partitions the graph into four quadrants. The vertex cover $C=(L\setminus Z)\cup(R\cap Z)$ (shaded orange) contains exactly the left endpoint of matches in the unreached half and the right endpoint of matches in the reached half. By definition of $Z$, no edge can cross from $L\cap Z$ to $R\setminus Z$, certifying that every edge is covered and $\abs{C}=\abs{M}$.}
\label{ch15:fig-konig-cover}
\end{figure}

Every edge in the bipartite graph is therefore covered, and the cover size matches the matching cardinality exactly: $\tau(\Gph) = \abs{C} = \abs{M} = \nu(\Gph)$.

\begin{labbox}
Run this on the running graph with maximum matching $M=\set{a2,b1,c4,d3}$. Every
left vertex is matched, so there is no free left vertex to seed the search and
$Z=\varnothing$. Then $C=(L\setminus\varnothing)\cup\varnothing = L=\set{a,b,c,d}$,
size $4$. Indeed $L$ itself is a cover of size $4=\nu$, and no smaller cover
exists because the matching already has $4$ disjoint edges. Now try the
non-maximum $M=\set{a1,b3,c4}$: the alternating search from free vertex $d$
reaches $3$ (via $d3$), then $b$ (via matched $b3$), then $1$ (via $b1$), then
$a$ (via matched $a1$), then $2$ (via $a2$) -- vertex $2$ is free, so the search
exposes an augmenting path and certifies that no cover of size $3$ exists.
\end{labbox}

Hall's theorem is the special case that asks when a \emph{perfect} matching of
the left side exists. For $S\subseteq L$ let $N(S)=\set{r\in R: \exists s\in
S,\,sr\in E}$ be its neighborhood.

\begin{theorem}[Hall's marriage theorem, 1935]
\label{ch15:thm-hall}
A bipartite graph has a matching saturating every vertex of $L$ if and only if
$\abs{N(S)}\ge\abs{S}$ for every $S\subseteq L$.
\end{theorem}

\begin{proof}
\emph{Necessity.} If $L$ is saturated, the $\abs{S}$ matched partners of $S$ are
distinct and lie in $N(S)$, so $\abs{N(S)}\ge\abs{S}$.

\emph{Sufficiency.} Suppose Hall's condition holds but the maximum matching
misses some left vertex, so $\nu<\abs{L}$. By K\H{o}nig there is a cover $C$ with
$\abs{C}=\nu<\abs{L}$. Split $C=C_L\cup C_R$ with $C_L=C\cap L$, $C_R=C\cap R$.
Let $S=L\setminus C_L$; then $\abs{S}=\abs{L}-\abs{C_L}>\abs{C_R}$. Every edge
from $S$ must be covered, and no vertex of $S$ is in $C$, so all such edges are
covered on the right by $C_R$; hence $N(S)\subseteq C_R$ and $\abs{N(S)}\le
\abs{C_R}<\abs{S}$, violating Hall. So $\nu=\abs{L}$ and $L$ is saturated.
\end{proof}

\begin{corollary}[Defect form / Ore]
\label{ch15:cor-defect}
The maximum matching saturates all but $\max_{S\subseteq L}\,\bigl(\abs{S}-
\abs{N(S)}\bigr)$ vertices of $L$; equivalently
$\nu=\abs{L}-\max_{S\subseteq L}(\abs{S}-\abs{N(S)})$.
\end{corollary}

Figure~\ref{ch15:fig-hall-violator} illustrates a Hall bottleneck: a subset $S\subseteq L$ of $3$ vertices has only $2$ collective neighbors in $R$, so by the pigeonhole principle no matching saturates $L$.

\begin{figure}[H]
\centering
\begin{tikzpicture}[>=Stealth, thick, scale=0.92,
  lnode/.style={circle, draw=nodeblue, fill=blue!12, minimum size=7.5mm, font=\small\bfseries},
  rnode/.style={circle, draw=targetgreen, fill=green!12, minimum size=7.5mm, font=\small\bfseries}]

  % Column headers
  \node[font=\footnotesize\bfseries, text=sourceblue] at (0, 3.8) {$L$};
  \node[font=\footnotesize\bfseries, text=targetgreen] at (4.5, 3.8) {$R$};

  % Subsets bounding boxes
  \draw[sourceblue, fill=blue!5, rounded corners=6pt, dashed, line width=1pt] (-0.5, 0.4) rectangle (0.8, 3.5);
  \node[sourceblue, font=\scriptsize\bfseries] at (-0.1, 3.2) {$S$};

  \draw[targetgreen, fill=green!5, rounded corners=6pt, dashed, line width=1pt] (3.7, 1.4) rectangle (5.3, 3.5);
  \node[targetgreen, font=\scriptsize\bfseries] at (4.9, 3.2) {$N(S)$};

  % Vertices in S
  \node[lnode] (u1) at (0.1, 2.6) {$u_1$};
  \node[lnode] (u2) at (0.1, 1.7) {$u_2$};
  \node[lnode] (u3) at (0.1, 0.8) {$u_3$};

  % Vertex outside S
  \node[lnode] (u4) at (0.1, -0.2) {$u_4$};

  % Vertices in N(S)
  \node[rnode] (v1) at (4.4, 2.6) {$v_1$};
  \node[rnode] (v2) at (4.4, 1.7) {$v_2$};

  % Vertex outside N(S)
  \node[rnode] (v3) at (4.4, -0.2) {$v_3$};

  % Edges from S (incident only to N(S))
  \draw[trapred, line width=2pt] (u1) -- node[midway, above=1pt, font=\tiny\bfseries, text=trapred] {$\in M$} (v1);
  \draw[edgegray] (u2) -- (v1);
  \draw[trapred, line width=2pt] (u2) -- node[midway, above=1pt, font=\tiny\bfseries, text=trapred] {$\in M$} (v2);
  \draw[edgegray] (u3) -- (v2);

  % Edge outside S
  \draw[trapred, line width=2pt] (u4) -- node[midway, above=1pt, font=\tiny\bfseries, text=trapred] {$\in M$} (v3);

  % Stranded callout on the right
  \node[anchor=west, draw=trapred, fill=red!6, rounded corners=4pt, inner sep=4pt, font=\scriptsize, text=gray!85!black, align=left] at (5.8, 1.2) {
    \textbf{\textcolor{trapred}{Hall Bottleneck:}}\\[2pt]
    $\bullet$ Subset size: $\abs{S} = 3$\\
    $\bullet$ Neighbors: $\abs{N(S)} = 2$\\
    $\bullet$ Deficiency: $\abs{S} - \abs{N(S)} = 1$\\[2pt]
    \textcolor{trapred}{\textbf{$u_3$ is stranded; no partner remains!}}
  };
\end{tikzpicture}
\caption{Hall's marriage condition violation. The subset $S=\{u_1,u_2,u_3\}\subseteq L$ has collective neighborhood $N(S)=\{v_1,v_2\}\subseteq R$. Because $\abs{N(S)}=2 < 3=\abs{S}$, the pigeonhole principle prevents any matching from saturating $S$. At least one vertex ($u_3$) is left permanently unmatched, demonstrating the defect bound $\nu \le \abs{L} - (\abs{S}-\abs{N(S)})$.}
\label{ch15:fig-hall-violator}
\end{figure}

Because no matching can map $3$ vertices into $2$ distinct targets, the defect is at least $\abs{S}-\abs{N(S)}=3-2=1$, directly certifying that the maximum matching size cannot exceed $3$.

\begin{connectionbox}
K\H{o}nig's theorem is the totally-unimodular special case of LP duality: the
matching polytope of a bipartite graph has integral vertices, so the LP relaxation
of maximum matching and its dual (fractional vertex cover) both attain integer
optima, and strong duality gives $\nu=\tau$. Total unimodularity in the companion
linear-algebra volume and the integral-flow theorem in Chapter~12 express the same
underlying integrality property. We make the flow connection explicit in
Section~\ref{ch15:flow}.
\end{connectionbox}

\section{The augmenting-path algorithm}
\label{ch15:kuhn}

Berge converts maximum matching into ``repeatedly find an augmenting path.''
In a bipartite graph an augmenting path is easy to find with a single search
rooted at a free left vertex, because of the orientation trick: walk forward
along unmatched edges $L\to R$ and backward along matched edges $R\to L$. If such
a walk reaches a free right vertex, we have an augmenting path.

The classic implementation (Kuhn's algorithm, the bipartite specialization of
the Hungarian method's first phase) does a depth-first search from each left
vertex in turn, trying to assign it a partner, recursively bumping current
occupants.

\begin{algorithm}[H]
\caption{\textsc{Augment}$(u)$ -- try to match left vertex $u$ by DFS}
\label{ch15:alg-augment}
\DontPrintSemicolon
\KwIn{left vertex $u$; arrays $\mathit{matchR}[\,]$ (current partner of each right
vertex, or $\bot$) and $\mathit{seen}[\,]$ (right vertices visited this search).}
\KwOut{\textsc{true} if $u$ can be matched, updating $\mathit{matchR}$.}
\ForEach{$r\in N(u)$}{
  \If{not $\mathit{seen}[r]$}{
    $\mathit{seen}[r]\gets\textsc{true}$\;
    \If{$\mathit{matchR}[r]=\bot$ \textbf{or} \textsc{Augment}$(\mathit{matchR}[r])$}{
      $\mathit{matchR}[r]\gets u$\;
      \Return \textsc{true}\;
    }
  }
}
\Return \textsc{false}\;
\end{algorithm}

\begin{algorithm}[H]
\caption{\textsc{MaxBipartiteMatching}$(L,R,E)$}
\label{ch15:alg-kuhn}
\DontPrintSemicolon
$\mathit{matchR}[r]\gets\bot$ for all $r\in R$\;
$\mathit{size}\gets 0$\;
\ForEach{$u\in L$}{
  reset $\mathit{seen}[r]\gets\textsc{false}$ for all $r\in R$\;
  \If{\textsc{Augment}$(u)$}{$\mathit{size}\gets\mathit{size}+1$}
}
\Return $\mathit{size}$ and $\mathit{matchR}$\;
\end{algorithm}

\begin{whybox}
Why does the recursive bump find an augmenting path? When \textsc{Augment}$(u)$
visits an unmatched neighbor $r$ it grabs it -- an augmenting path of length $1$.
When it visits a \emph{matched} neighbor $r$ whose current owner is $w$, it asks
$w$ to relocate; if $w$ succeeds, the recursion has built a longer alternating
path $u\!-\!r\!-\!w\!-\!\dots\!-\!r'$ ending at a free right vertex $r'$, and
unwinding the recursion flips exactly that path. The $\mathit{seen}$ array stops
the DFS from revisiting a right vertex, which both guarantees termination and
keeps each search linear.
\end{whybox}

\begin{whybox}[Correctness of the outer loop]
The outer loop processes left vertices one at a time and never \emph{unmatches}
anyone permanently: a successful \textsc{Augment} only relocates owners along an
augmenting path, so the matching size strictly increases and previously matched
left vertices stay matched. The invariant is that after processing the first $k$
left vertices, the current matching is maximum on the subgraph induced by those
$k$ vertices and all of $R$. A failed \textsc{Augment}$(u)$ certifies (by the
same alternating-reachability argument as K\H{o}nig) that no augmenting path from
$u$ exists relative to the current matching, and adding later left vertices can
never create one rooted at the already-failed $u$. So one pass suffices and the
final matching is maximum.
\end{whybox}

\paragraph{Trace on the running graph.}
Adjacency in listed order: $a:\set{1,2}$, $b:\set{1,3}$, $c:\set{3,4}$,
$d:\set{3}$. Process $a,b,c,d$.

\begin{center}
\begin{tabular}{@{}llll@{}}
\toprule
step & call & what happens & $\mathit{matchR}$ after \\
\midrule
1 & \textsc{Augment}$(a)$ & $1$ free $\to$ take it & $1\!:\!a$ \\
2 & \textsc{Augment}$(b)$ & $1$ owned by $a$; bump $a$: $a$ takes $2$ (free) & $1\!:\!b,\;2\!:\!a$ \\
3 & \textsc{Augment}$(c)$ & $3$ free $\to$ take it & $1\!:\!b,2\!:\!a,3\!:\!c$ \\
4 & \textsc{Augment}$(d)$ & $3$ owned by $c$; bump $c$: $c$ takes $4$ (free) & $1\!:\!b,2\!:\!a,3\!:\!d,4\!:\!c$ \\
\bottomrule
\end{tabular}
\end{center}

Step 2 is the augmenting path $b\!-\!1\!-\!a\!-\!2$ (length $3$); step 4 is
$d\!-\!3\!-\!c\!-\!4$. The final matching $\set{b1,a2,d3,c4}$ is perfect, size
$4$ -- exactly the optimum greedy missed.

\begin{complexitybox}
Each call to \textsc{MaxBipartiteMatching} runs $\abs{L}$ outer iterations. Each
iteration is a DFS that, thanks to $\mathit{seen}$, touches every edge at most
once: $\Oh(m)$. Total \[ \Oh(\abs{L}\cdot m)=\Oh(V\!E). \] With $V=n$ and
$m\le\frac{n^2}{4}$ this is $\Oh(n^3)$ in the dense worst case. Space is $\Oh(V)$
beyond the graph. The bound is tight for adversarial vertex orderings: paths can
grow to length $\Th(V)$ and you may pay $\Th(V)$ per augmentation.
\end{complexitybox}

\begin{trapbox}
Resetting $\mathit{seen}$ \emph{between} outer iterations (once per left vertex)
is correct and necessary; resetting it \emph{inside} the recursion, or forgetting
to reset it, both break the algorithm -- the first reintroduces exponential
blowup, the second makes \textsc{Augment} return spurious failures. Also note
\textsc{Augment} marks $r$ as seen \emph{before} recursing, never after; marking
after recursion lets two branches chase the same right vertex.
\end{trapbox}

\section{Hopcroft--Karp: augment in batches}
\label{ch15:hk}

Kuhn pays $\Oh(m)$ per single augmenting path and needs up to $\nu$ of them.
Hopcroft and Karp (1973) observed that we are wasteful: instead of finding
\emph{one} shortest augmenting path per phase, find a \emph{maximal set of
vertex-disjoint shortest augmenting paths} all at once, and flip them together.
Two structural facts make this both correct and fast: the shortest augmenting
path length strictly increases from phase to phase, and after $\Oh(\sqrt{V})$
phases the matching is maximum.

\begin{intuitionbox}
Think of it as flood-then-drain. A breadth-first sweep from all free left
vertices, alternating unmatched/matched as it descends, lays down a layered
``distance'' structure and finds the shortest distance $\ell$ to any free right
vertex. Then a depth-first sweep greedily extracts as many length-$\ell$
augmenting paths as it can without reusing a vertex, and flips them all. Each BFS
sweep raises the shortest-path length, so the lengths go $1,3,5,\dots$ and cannot
go on long. Short paths are cheap and plentiful; the expensive long paths are
provably few.
\end{intuitionbox}

\begin{algorithm}[H]
\caption{\textsc{HopcroftKarp}$(L,R,E)$}
\label{ch15:alg-hk}
\DontPrintSemicolon
$\mathit{matchL}[u]\gets\bot,\ \mathit{matchR}[r]\gets\bot$ for all $u,r$\;
$\mathit{size}\gets 0$\;
\While{\textsc{BFS}() reports a free-to-free layered path exists}{
  \ForEach{free $u\in L$}{
    \If{\textsc{DFS}$(u)$}{$\mathit{size}\gets\mathit{size}+1$}
  }
}
\Return $\mathit{size}$\;
\end{algorithm}

\begin{algorithm}[H]
\caption{\textsc{BFS}() -- build layers, return whether any augmenting path exists}
\label{ch15:alg-hk-bfs}
\DontPrintSemicolon
queue $Q\gets$ all free $u\in L$; set $\mathit{dist}[u]\gets 0$ for these,
$\mathit{dist}[\cdot]\gets\infty$ otherwise (including sentinel $\mathit{dist}[\bot]\gets\infty$)\;
\While{$Q$ nonempty}{
  pop $u$\;
  \If{$\mathit{dist}[u] < \mathit{dist}[\bot]$}{
    \ForEach{$r\in N(u)$}{
      $w\gets\mathit{matchR}[r]$\;
      \If{$\mathit{dist}[w]=\infty$}{
        $\mathit{dist}[w]\gets\mathit{dist}[u]+1$\;
        \If{$w\neq\bot$}{push $w$ onto $Q$}
      }
    }
  }
}
\Return $\mathit{dist}[\bot]\neq\infty$\;
\end{algorithm}

\begin{algorithm}[H]
\caption{\textsc{DFS}$(u)$ -- extract one shortest augmenting path from $u$}
\label{ch15:alg-hk-dfs}
\DontPrintSemicolon
\If{$u=\bot$}{\Return \textsc{true}}
\ForEach{$r\in N(u)$}{
  $w\gets\mathit{matchR}[r]$\;
  \If{$\mathit{dist}[w]=\mathit{dist}[u]+1$ \textbf{and} \textsc{DFS}$(w)$}{
    $\mathit{matchR}[r]\gets u$;\ $\mathit{matchL}[u]\gets r$;\ \Return
    \textsc{true}\;
  }
}
$\mathit{dist}[u]\gets\infty$;\ \Return \textsc{false}\;
\end{algorithm}

The BFS computes $\mathit{dist}[u]$ = number of $L\to R\to L$ hops along a
shortest alternating path from a free left vertex to $u$, recording the first
free-right-vertex layer in the sentinel $\mathit{dist}[\bot]$ so longer paths are
never mixed into the same phase. The DFS then steps strictly from $\mathit{dist}[u]$
to $\mathit{dist}[u]+1$ all the way to $\bot$, extracting only shortest augmenting
paths of edge length $2\,\mathit{dist}[\bot]-1$; setting $\mathit{dist}[u]\gets\infty$ on failure
prunes dead ends so no edge is explored twice in a phase. Figure~\ref{ch15:fig-hopcroft-karp}
illustrates this two-phase structure: the BFS lays down the level DAG, and the DFS
greedily extracts a maximal vertex-disjoint family of augmenting paths in a single pass.
A phase therefore runs in $\Oh(m)$.

\begin{figure}[H]
\centering
\begin{tikzpicture}[
  every node/.style={font=\small},
  L/.style={circle,draw=sourceblue,thick,fill=blue!10,minimum size=6.5mm},
  R/.style={circle,draw=targetgreen,thick,fill=green!10,minimum size=6.5mm},
  fL/.style={circle,draw=sourceblue,very thick,fill=blue!22,minimum size=6.5mm},
  fR/.style={circle,draw=targetgreen,very thick,fill=green!22,minimum size=6.5mm},
  unm/.style={-{Stealth[length=2mm]},edgegray,thick},
  mat/.style={-{Stealth[length=2mm]},trapred,line width=1.5pt},
  pA/.style={-{Stealth[length=2.5mm]},sourceblue,line width=1.8pt},
  pB/.style={-{Stealth[length=2.5mm]},quotientorange,line width=1.8pt},
  lane/.style={draw=black!15,fill=black!2,rounded corners=4pt}
]
% Panel (a)
\begin{scope}
  \fill[lane] (-0.6,-0.6) rectangle (0.6,2.6);
  \fill[lane] (1.2,-0.6) rectangle (2.4,2.6);
  \fill[lane] (3.0,-0.6) rectangle (4.2,2.6);
  \fill[lane] (4.8,-0.6) rectangle (6.0,2.6);

  \node[font=\footnotesize\bfseries,align=center] at (0,2.9) {Layer 0\\[-1pt]\color{sourceblue}$L_{\mathrm{free}}$};
  \node[font=\footnotesize\bfseries,align=center] at (1.8,2.9) {Layer 1\\[-1pt]\color{targetgreen}$R$};
  \node[font=\footnotesize\bfseries,align=center] at (3.6,2.9) {Layer 2\\[-1pt]\color{sourceblue}$L_{\mathrm{matched}}$};
  \node[font=\footnotesize\bfseries,align=center] at (5.4,2.9) {Layer 3\\[-1pt]\color{targetgreen}$R_{\mathrm{free}}$};

  \node[fL] (u1) at (0,2.0) {$u_1$};
  \node[fL] (u2) at (0,0.5) {$u_2$};

  \node[R] (v1) at (1.8,2.0) {$v_1$};
  \node[R] (v2) at (1.8,0.5) {$v_2$};

  \node[L] (w1) at (3.6,2.0) {$w_1$};
  \node[L] (w2) at (3.6,0.5) {$w_2$};

  \node[fR] (z1) at (5.4,2.0) {$z_1$};
  \node[fR] (z2) at (5.4,0.5) {$z_2$};

  \draw[unm] (u1) -- (v1);
  \draw[unm] (u1) -- (v2);
  \draw[unm] (u2) -- (v2);

  \draw[mat] (v1) -- (w1) node[midway,above=1pt,font=\scriptsize,trapred] {$\in M$};
  \draw[mat] (v2) -- (w2) node[midway,below=1pt,font=\scriptsize,trapred] {$\in M$};

  \draw[unm] (w1) -- (z1);
  \draw[unm] (w2) -- (z1);
  \draw[unm] (w2) -- (z2);

  \node[font=\footnotesize,align=center] at (2.7,-1.2) {\textbf{(a)} BFS layered DAG: alternating paths\\from free left vertices to free right targets};
\end{scope}

% Panel (b)
\begin{scope}[xshift=7.4cm]
  \fill[lane] (-0.6,-0.6) rectangle (0.6,2.6);
  \fill[lane] (1.2,-0.6) rectangle (2.4,2.6);
  \fill[lane] (3.0,-0.6) rectangle (4.2,2.6);
  \fill[lane] (4.8,-0.6) rectangle (6.0,2.6);

  \node[font=\footnotesize\bfseries,align=center] at (0,2.9) {Layer 0\\[-1pt]\color{sourceblue}$L_{\mathrm{free}}$};
  \node[font=\footnotesize\bfseries,align=center] at (1.8,2.9) {Layer 1\\[-1pt]\color{targetgreen}$R$};
  \node[font=\footnotesize\bfseries,align=center] at (3.6,2.9) {Layer 2\\[-1pt]\color{sourceblue}$L_{\mathrm{matched}}$};
  \node[font=\footnotesize\bfseries,align=center] at (5.4,2.9) {Layer 3\\[-1pt]\color{targetgreen}$R_{\mathrm{free}}$};

  \node[fL] (bu1) at (0,2.0) {$u_1$};
  \node[fL] (bu2) at (0,0.5) {$u_2$};

  \node[R] (bv1) at (1.8,2.0) {$v_1$};
  \node[R] (bv2) at (1.8,0.5) {$v_2$};

  \node[L] (bw1) at (3.6,2.0) {$w_1$};
  \node[L] (bw2) at (3.6,0.5) {$w_2$};

  \node[fR] (bz1) at (5.4,2.0) {$z_1$};
  \node[fR] (bz2) at (5.4,0.5) {$z_2$};

  % Background unselected edges
  \draw[edgegray!40,thick] (bu1) -- (bv2);
  \draw[edgegray!40,thick] (bw2) -- (bz1);

  % Extracted Path 1 (P_1) in blue
  \draw[pA] (bu1) -- (bv1);
  \draw[pA] (bv1) -- (bw1);
  \draw[pA] (bw1) -- (bz1);

  % Extracted Path 2 (P_2) in orange
  \draw[pB] (bu2) -- (bv2);
  \draw[pB] (bv2) -- (bw2);
  \draw[pB] (bw2) -- (bz2);

  \node[sourceblue,font=\footnotesize\bfseries,above=1pt] at (2.7,2.2) {Augmenting Path $P_1$};
  \node[quotientorange,font=\footnotesize\bfseries,below=1pt] at (2.7,0.3) {Augmenting Path $P_2$};

  \node[font=\footnotesize,align=center] at (2.7,-1.2) {\textbf{(b)} DFS extraction: maximal vertex-disjoint\\paths $P_1, P_2$ flipped simultaneously};
\end{scope}
\end{tikzpicture}
\caption{The two phases of Hopcroft--Karp. (a)~A forward breadth-first
search starts at free left vertices $L_{\mathrm{free}}$ (Layer 0), traversing
unmatched edges $L \to R$ and matched edges $R \to L$ to build a layered DAG of
shortest alternating distances until reaching free right vertices
$R_{\mathrm{free}}$ (Layer 3). (b)~A depth-first search then traverses only DAG
edges with $\mathit{dist}[w] = \mathit{dist}[u]+1$, greedily collecting a maximal
family of vertex-disjoint shortest augmenting paths ($P_1$ in blue, $P_2$ in orange)
to augment simultaneously in $\Oh(m)$ time.}
\label{ch15:fig-hopcroft-karp}
\end{figure}

\paragraph{Trace on the running graph.}
Start $M=\varnothing$; all of $a,b,c,d$ free.

\emph{Phase 1.} BFS: every free left vertex has $\mathit{dist}=0$; all their
right neighbors are unmatched, so the shortest augmenting length is $1$
($\mathit{found}=\textsc{true}$). DFS in order $a,b,c,d$, each grabbing a
disjoint free right neighbor. $a$ takes $1$; $b$ takes $3$; $c$'s neighbor $3$ is
used this phase, so $c$ takes $4$; $d$'s only neighbor $3$ is used this phase, so
$d$ fails. After phase 1: $M=\set{a1,b3,c4}$, size $3$, $d$ still free.

\emph{Phase 2.} Only free left vertex is $d$ ($\mathit{dist}[d]=0$). BFS from $d$:
neighbor $3$ is matched to $b$, set $\mathit{dist}[b]=1$, push $b$. From $b$:
neighbor $1$ matched to $a$, set $\mathit{dist}[a]=2$, push $a$. From $a$:
neighbor $2$ unmatched $\Rightarrow$ $\mathit{found}=\textsc{true}$ (shortest
length now $5$). DFS from $d$: $d\!\to\!3$ (owner $b$,
$\mathit{dist}[b]=1=\mathit{dist}[d]+1$), recurse $b\!\to\!1$ (owner $a$,
$\mathit{dist}[a]=2$), recurse $a\!\to\!2$ (unmatched) -- success. Flip the chain
$d\!-\!3\!-\!b\!-\!1\!-\!a\!-\!2$: now $d3,b1,a2$ matched, $c4$ untouched.
$M=\set{a2,b1,c4,d3}$, size $4$.

\emph{Phase 3.} No free left vertex remains; BFS returns \textsc{false}; stop.
Two phases of work, as the $\sqrt{V}$ bound predicts ($\sqrt{4}=2$).

\subsection{Why the phase count is \texorpdfstring{$\Oh(\sqrt{V})$}{O(sqrt V)}}
\label{ch15:hk-analysis}

The analysis rests on a monotonicity lemma and a counting argument. Let
$\ell_i$ be the length (number of edges) of the shortest augmenting path at the
start of phase $i$.

\begin{lemma}[Lengths strictly increase]
\label{ch15:lem-monotone}
If phase $i$ augments along a maximal set of vertex-disjoint shortest augmenting
paths of length $\ell_i$, then the shortest augmenting path of the new matching
has length $\ell_{i+1}>\ell_i$.
\end{lemma}

\begin{proof}
Let $M$ be the matching before phase $i$, $\mathcal{P}$ the maximal family of
$k=\abs{\mathcal{P}}$ vertex-disjoint shortest $M$-augmenting paths (each of length
$\ell_i$), and $M' = M \oplus \bigl(\bigcup_{P\in\mathcal{P}} P\bigr)$. Take any
$M'$-augmenting path $Q$, and set $M'' = M' \oplus Q$, so $\abs{M''} = \abs{M} + k + 1$.
In the symmetric difference $M \oplus M'' = \bigl(\bigcup_{P\in\mathcal{P}} P\bigr) \oplus Q$,
every component is an alternating path or even cycle, and because $\abs{M''}-\abs{M}=k+1$,
at least $k+1$ components are vertex-disjoint $M$-augmenting paths $P'_1,\dots,P'_{k+1}$.
Every $M$-augmenting path has length at least $\ell_i$, so $\sum_{j=1}^{k+1}\abs{P'_j} \ge (k+1)\ell_i$.

Now compare that lower bound against the edge count of $\bigl(\bigcup_{P\in\mathcal{P}} P\bigr) \oplus Q$:
\begin{itemize}
\item If $Q$ shares \emph{no} vertex with $\bigcup_{P\in\mathcal{P}} P$, then $Q$ uses no edge touched
by phase $i$, so $Q$ was already an $M$-augmenting path vertex-disjoint from all of
$\mathcal{P}$. By maximality of $\mathcal{P}$ among length-$\ell_i$ paths, $\abs{Q} > \ell_i$.
\item Otherwise $Q$ shares at least one vertex $x$ with some $P\in\mathcal{P}$. Because $P$ is
saturated by $M'$, $x$ is incident to a unique edge $e\in M'$, which lies on $P$; and since $x$
cannot be an $M'$-free endpoint of $Q$, the $M'$-alternating path $Q$ must also traverse $e$.
Thus $e \in \bigl(\bigcup_{P\in\mathcal{P}} P\bigr) \cap Q$, canceling from both sides of the
symmetric difference:
\[
  (k+1)\ell_i \;\le\; \sum_{j=1}^{k+1}\abs{P'_j} \;\le\; \Bigl\lvert\Bigl(\bigcup_{P\in\mathcal{P}} P\Bigr)\oplus Q\Bigr\rvert
  \;=\; k\ell_i + \abs{Q} - 2\Bigl\lvert\Bigl(\bigcup_{P\in\mathcal{P}} P\Bigr)\cap Q\Bigr\rvert
  \;\le\; k\ell_i + \abs{Q} - 2,
\]
which rearranges immediately to $\abs{Q} \ge \ell_i + 2 > \ell_i$.
\end{itemize}
In both cases $\ell_{i+1} > \ell_i$.
\end{proof}

Augmenting paths have odd length, so the lengths run $1,3,5,\dots$, strictly
increasing.

\begin{theorem}[Hopcroft--Karp phase bound]
\label{ch15:thm-hk-phases}
The algorithm terminates after at most $\Oh(\sqrt{\nu})=\Oh(\sqrt{V})$ phases,
where $\nu$ is the size of the maximum matching.
\end{theorem}

\begin{proof}
Run the algorithm until the shortest augmenting path first has length exceeding
$\sqrt{\nu}$. By Lemma~\ref{ch15:lem-monotone} the lengths are
$\ge 1,3,5,\dots$, so this happens within $\lceil\sqrt{\nu}\,\rceil$ phases. Let
$M$ be the current matching at that moment and $M^\ast$ a maximum matching. As in
Berge's proof, $M\oplus M^\ast$ is a disjoint union of paths and even cycles and
contains exactly $\nu-\abs{M}$ vertex-disjoint $M$-augmenting paths. Each such
path now has length $>\sqrt{\nu}$, hence uses more than $\tfrac12\sqrt{\nu}$
edges of $M$. These paths are vertex-disjoint and all their $M$-edges lie in the
$\le\nu$ edges of $M$, so the number of paths is at most
$\nu/(\tfrac12\sqrt{\nu})=2\sqrt{\nu}$. Thus $\nu-\abs{M}\le 2\sqrt{\nu}$: at most
$2\sqrt{\nu}$ augmentations remain, and each later phase performs at least one,
so at most $2\sqrt{\nu}$ further phases finish. Total phases
$\le\lceil\sqrt{\nu}\,\rceil+2\sqrt{\nu}=\Oh(\sqrt{\nu})=\Oh(\sqrt{V})$.
\end{proof}

\begin{complexitybox}
Each phase is one BFS plus one DFS sweep, each $\Oh(m)$ amortized (the
$\mathit{dist}\gets\infty$ pruning ensures every edge is examined a constant
number of times per phase). With $\Oh(\sqrt{V})$ phases,
\[
\boxed{\;\Oh\!\bigl(E\sqrt{V}\bigr)\;}
\]
time and $\Oh(V)$ extra space. For dense graphs this is $\Oh(V^{2.5})$ versus
Kuhn's $\Oh(V^3)$. On bounded-degree or random graphs the constant is small and
Hopcroft--Karp is the practical default.
\end{complexitybox}

\begin{historybox}
Bipartite matching grew out of both matrix theory and combinatorics.
In 1931, D{\'e}nes K{\H{o}}nig established that the maximum size of a matching equals the minimum size
of a vertex cover, motivated by Frobenius's work on matrix determinants. In 1935, Philip Hall
formulated his marriage theorem for systems of distinct representatives, providing the
neighborhood condition for when every element of one bipartition can be matched. Claude Berge in 1957
established the algorithmic viewpoint by proving that a matching is maximum if and only if
it admits no augmenting path.

In 1973, John Hopcroft and Richard Karp proved that searching for maximal families of shortest
augmenting paths simultaneously in breadth-first layers bounds the number of phases by $\Oh(\sqrt{V})$,
achieving an overall running time of $\Oh(E\sqrt{V})$. In the Soviet Union, Alexander Karzanov
independently obtained the same bound using unit-capacity blocking flows, reflecting the direct
equivalence between Hopcroft--Karp and Dinitz's layered search. For five decades, $\Oh(E\sqrt{V})$
remained the best known general bound. In 2022, Li Chen, Rasmus Kyng, Yang Liu, Richard Peng,
Maximilian Probst Gutenberg, and Sushant Sachdeva improved the asymptotic bound via their
$m^{1+o(1)}$ almost-linear flow framework, though Hopcroft--Karp remains the practical implementation
of choice on real-world graphs.
\end{historybox}

\section{Matching via maximum flow}
\label{ch15:flow}

Every result above has a direct counterpart in network flow, and making the reduction explicit both
gives an alternative algorithm and explains \emph{why} the min-max theorems hold.

\paragraph{The construction.}
Build a directed flow network $\Gph'$:
\begin{itemize}
\item a source $s$ and sink $t$;
\item an arc $s\to u$ of capacity $1$ for every $u\in L$;
\item an arc $u\to r$ of capacity $1$ (or $\infty$) for every edge $ur\in E$,
directed $L\to R$;
\item an arc $r\to t$ of capacity $1$ for every $r\in R$.
\end{itemize}

\begin{figure}[H]
\centering
\begin{tikzpicture}[
  every node/.style={font=\small},
  S/.style={circle,draw=quotientorange,thick,fill=orange!12,minimum size=8mm},
  L/.style={circle,draw=sourceblue,thick,fill=blue!8,minimum size=7mm},
  R/.style={circle,draw=targetgreen,thick,fill=green!8,minimum size=7mm},
  a/.style={-{Stealth[length=2mm]},edgegray}
]
\node[S] (s) at (-2.3,1.5) {$s$};
\node[S] (t) at (5.5,1.5) {$t$};
\foreach \v/\y in {a/3,b/2,c/1,d/0} \node[L] (\v) at (0.3,\y) {$\v$};
\foreach \v/\y in {1/3,2/2,3/1,4/0} \node[R] (r\v) at (3.0,\y) {$\v$};
\foreach \v in {a,b,c,d} \draw[a] (s)--(\v);
\foreach \v in {1,2,3,4} \draw[a] (r\v)--(t);
\draw[a] (a)--(r1); \draw[a] (a)--(r2);
\draw[a] (b)--(r1); \draw[a] (b)--(r3);
\draw[a] (c)--(r3); \draw[a] (c)--(r4);
\draw[a] (d)--(r3);
\node at (-1.1,2.6) {\scriptsize cap $1$};
\node at (4.4,2.6) {\scriptsize cap $1$};
\end{tikzpicture}
\caption{The running graph turned into a unit-capacity $s$--$t$ network. Every
$s$--$t$ path has the form $s\to u\to r\to t$ and uses three unit-capacity arcs.}
\label{ch15:fig-flow}
\end{figure}

\begin{theorem}[Matching $=$ integral max flow]
\label{ch15:thm-flow}
Integral $s$--$t$ flows of value $k$ in $\Gph'$ correspond exactly to matchings
of size $k$ in $\Gph$; in particular the maximum flow value equals $\nu(\Gph)$.
\end{theorem}

\begin{proof}
All capacities are integral, so by the integral flow theorem a maximum flow can
be taken integral, with every arc carrying $0$ or $1$. Given such a flow $f$, let
$M=\set{ur\in E: f(u\to r)=1}$. Conservation at $u\in L$, whose only in-arc is the
capacity-$1$ arc $s\to u$, forces at most one out-arc of $u$ to carry flow, so
$u$ meets at most one edge of $M$; symmetrically each $r\in R$ feeds the single
capacity-$1$ arc $r\to t$, so $r$ meets at most one edge of $M$. Thus $M$ is a
matching, and $\abs{M}=\sum_u f(s\to u)=\abs{f}$. Conversely a matching $M$ gives
an integral flow of value $\abs{M}$ by routing one unit $s\to u\to r\to t$ for
each $ur\in M$. The correspondence is value-preserving, so $\max\abs{f}=\nu$.
\end{proof}

\begin{whybox}
This is why K\H{o}nig is ``just'' max-flow/min-cut. Give the middle $L\to R$ arcs
capacity $\infty$, so a finite minimum $s$--$t$ cut never cuts one; it consists of
cut $s\to u$ arcs and cut $r\to t$ arcs only. The corresponding left vertices and
right vertices together form a vertex cover, and min-cut $=$ max-flow becomes
$\tau=\nu$. Hall is the ``max flow saturates all of $L$'' specialization. The flow
viewpoint hands you the dual certificate for free.
\end{whybox}

\begin{whybox}[Augmenting paths are augmenting paths]
A residual $s$--$t$ path in $\Gph'$ alternates: forward $s\to u$, forward $u\to r$
(an unused matching edge), backward $r\to u'$ (a used edge, i.e. matched,
traversed against its flow), forward $u'\to r'$, $\dots$, forward $r''\to t$.
Strip the $s$ and $t$ endpoints and this is precisely an $M$-augmenting path in
$\Gph$: forward arcs are non-matching edges, backward arcs are matching edges,
and the endpoints $u,r''$ are free. Ford--Fulkerson on $\Gph'$ \emph{is} Berge's
algorithm; BFS augmenting (Edmonds--Karp) on unit capacities reproduces
Hopcroft--Karp's phase structure (Dinic blocking flow). The matching algorithms
and the flow algorithms are the same algorithms.
\end{whybox}

\begin{complexitybox}
Plain Ford--Fulkerson with BFS on this unit-capacity network gives the same
$\Oh(VE)$ as Kuhn. Dinic's blocking-flow algorithm on a unit-capacity network
runs in $\Oh(E\sqrt{V})$ -- identical to Hopcroft--Karp, and for the same reason
(the $\sqrt{V}$-phase bound is a theorem about unit-capacity blocking flows). Thus
reducing bipartite matching to unit-capacity flow runs the same augmenting phases with extra source and
sink vertices. The flow framing wins when the model needs capacities $>1$
(e.g. one job can absorb $b$ workers), which a plain matching cannot express.
\end{complexitybox}

\begin{applicationbox}
The flow reduction makes bipartite matching a versatile modeling tool.
Assigning radio frequencies to towers, scheduling exams to rooms and time slots,
routing virtual machines to physical hosts, recovering a perfect matching in a
Latin-square completion, and the ``system of distinct representatives'' problems
all become a single max-flow call once phrased as left/right compatibility.
Capacitated variants -- a hospital that can absorb $q$ residents, a warehouse
that ships $c$ units -- are immediate by raising the $s\to u$ or $r\to t$
capacities, and the integral-flow theorem keeps the answer whole. The
flow-modeling chapter of this book gives the full catalogue.
\end{applicationbox}

\begin{connectionbox}
The matching polytope of a bipartite graph is the set of fractional matchings
$\set{x\ge 0:\ \sum_{e\ni v}x_e\le 1}$, and its constraint matrix is the
incidence matrix of a bipartite graph -- which is \emph{totally unimodular}. Total
unimodularity forces every vertex of the polytope to be integral, so the LP and
the combinatorial problem coincide. This is a natural entry point to the
polyhedral and spectral material of the book: the same incidence matrix, its
associated Laplacian $\Lap=BB^{\!\top}$ for incidence matrix $B$, and the rank
facts that drive total unimodularity all reappear in the chapter on the graph
Laplacian. Matching is one of the clearest places where combinatorial optimization
and linear algebra meet.
\end{connectionbox}

\section{Putting it together}
\label{ch15:summary}

\begin{center}
\begin{tabular}{@{}lll@{}}
\toprule
algorithm & time & idea \\
\midrule
Kuhn / augmenting DFS & $\Oh(VE)$ & one augmenting path per left vertex \\
Hopcroft--Karp & $\Oh(E\sqrt{V})$ & maximal disjoint shortest paths per phase \\
Dinic on unit network & $\Oh(E\sqrt{V})$ & blocking flow $=$ Hopcroft--Karp \\
\bottomrule
\end{tabular}
\end{center}

The central result is Berge's theorem
(Theorem~\ref{ch15:thm-berge}): a matching is maximum iff no augmenting path
survives. K\H{o}nig (Theorem~\ref{ch15:thm-konig}) and Hall
(Theorem~\ref{ch15:thm-hall}) give the dual obstruction -- a minimum vertex cover
of equal size, and the neighborhood condition for saturating one side. The
algorithms differ only in how they search for augmenting paths, and the flow
reduction (Theorem~\ref{ch15:thm-flow}) shows they all execute the same underlying
residual-graph search.

\begin{recoverybox}
Before the exercises, re-derive the fact the whole chapter depends on: take any two
matchings, look at their symmetric difference, and convince yourself it is a
disjoint union of paths and even cycles. Every theorem in this chapter -- Berge,
K\H{o}nig (via the alternating-reachability set $Z$), Hall (via K\H{o}nig), and
the Hopcroft--Karp phase bound -- comes from this picture with a different counting
argument.
\end{recoverybox}

Every matching in this chapter treated edges identically, aiming only to pair as many vertices as possible.
In practical assignments, compatibility comes in degrees: applicants have test scores, drivers have passenger ratings, and workers have wage expectations.
The next chapter assigns each edge a numerical weight, introducing the Hungarian algorithm and primal--dual methods to find maximum-weight bipartite matchings.

\section{Exercises}
\label{ch15:exercises}

\begin{exercise}[Greedy is a 2-approximation]
\label{ch15:ex-greedy}
A \emph{maximal} matching is one you cannot enlarge by adding a single edge.
Prove that any maximal matching has at least $\nu/2$ edges, and exhibit a
bipartite graph where greedy returns exactly $\nu/2$. (Hint: each greedy edge can
block at most two optimal edges.)
\end{exercise}

\begin{exercise}[Implement and instrument Kuhn]
\label{ch15:ex-kuhn}
Implement Algorithm~\ref{ch15:alg-kuhn}. Add a counter of edges examined and run
it on a graph designed to force $\Th(VE)$ work. Confirm empirically that
reordering $L$ by ascending degree usually slashes the runtime, even though the
worst case is unchanged.
\end{exercise}

\begin{exercise}[Break a buggy \textsc{seen} reset]
\label{ch15:ex-seen}
Modify \textsc{Augment} (Algorithm~\ref{ch15:alg-augment}) to reset
$\mathit{seen}[r]$ to \textsc{false} just before returning \textsc{true}. Find a
small graph on which the modified routine reports a matching smaller than
maximum, and explain which augmenting path it fails to find.
\end{exercise}

\begin{exercise}[K\H{o}nig from a run]
\label{ch15:ex-konig}
After computing a maximum matching with any algorithm, give an $\Oh(V+E)$
procedure that outputs a minimum vertex cover of equal size, following the
reachability set $Z$ in the proof of Theorem~\ref{ch15:thm-konig}. Run it on the
chapter's graph and verify you recover a cover of size $4$.
\end{exercise}

\begin{exercise}[Hall, tightly]
\label{ch15:ex-hall}
Construct a bipartite graph with $\abs{L}=\abs{R}=n$ in which exactly one subset
$S\subseteq L$ violates Hall's condition, with $\abs{N(S)}=\abs{S}-1$. What is the
maximum matching size, and which vertices does Corollary~\ref{ch15:cor-defect}
predict are left unsaturated?
\end{exercise}

\begin{exercise}[Phase lengths in Hopcroft--Karp]
\label{ch15:ex-hk-lengths}
Build a bipartite graph on which Hopcroft--Karp uses augmenting paths of lengths
$1,3,5$ across three distinct phases (so the bound is exercised, not just the
trivial single-phase case). Trace the BFS distance labels at the start of each
phase using Algorithm~\ref{ch15:alg-hk-bfs}.
\end{exercise}

\begin{exercise}[Tightness of $\Oh(E\sqrt{V})$]
\label{ch15:ex-hk-tight}
It is known that on some bipartite graphs Hopcroft--Karp truly needs
$\Th(\sqrt{V})$ phases. Sketch a ``staircase'' of nested long alternating paths and
argue that the shortest augmenting length grows by exactly $2$ each phase, so
$\Th(\sqrt{V})$ phases occur.
\end{exercise}

\begin{exercise}[Capacitated assignment via flow]
\label{ch15:ex-cap}
Suppose each right vertex $r$ can be matched to up to $q_r$ left vertices (a
\emph{$b$-matching} on the right). Show how to model this with the network of
Section~\ref{ch15:flow} by changing one set of capacities, and argue the
integral-flow theorem still returns an integer solution. What is the analogue of
K\H{o}nig's theorem here?
\end{exercise}

\begin{exercise}[Augmenting paths fail the easy way in general graphs]
\label{ch15:ex-blossom}
Give a non-bipartite graph (an odd cycle plus a pendant edge will do) and a
matching where the naive ``DFS alternating from a free vertex'' of
Algorithm~\ref{ch15:alg-augment} reports failure even though an augmenting path
exists. Explain why the odd cycle defeats the bipartite reachability argument --
this is the blossom that Chapter~17 contracts.
\end{exercise}

\begin{exercise}[Matching and the incidence matrix]
\label{ch15:ex-incidence}
Let $B$ be the $\abs{V}\times\abs{E}$ vertex-edge incidence matrix of a bipartite
graph. Prove $B$ is totally unimodular (every square submatrix has determinant in
$\set{-1,0,1}$), and deduce that the fractional matching LP has an integral
optimum. Relate this to detecting a perfect matching via the determinant of the
bipartite biadjacency matrix (the Edmonds matrix idea).
\end{exercise}

\begin{exercise}[Online matching, a glimpse ahead]
\label{ch15:ex-online}
Left vertices arrive one at a time and must be matched irrevocably on arrival, or
dropped. Show that the deterministic greedy rule is $2$-competitive and no
better, then look up and summarize the \textsc{Ranking} algorithm's $1-1/e$
guarantee. Why does the offline $\Oh(E\sqrt{V})$ machinery not apply?
\end{exercise}
